% Options for packages loaded elsewhere
\PassOptionsToPackage{unicode}{hyperref}
\PassOptionsToPackage{hyphens}{url}
%
\documentclass[
]{article}
\usepackage{amsmath,amssymb}
\usepackage{lmodern}
\usepackage{iftex}
\ifPDFTeX
  \usepackage[T1]{fontenc}
  \usepackage[utf8]{inputenc}
  \usepackage{textcomp} % provide euro and other symbols
\else % if luatex or xetex
  \usepackage{unicode-math}
  \defaultfontfeatures{Scale=MatchLowercase}
  \defaultfontfeatures[\rmfamily]{Ligatures=TeX,Scale=1}
\fi
% Use upquote if available, for straight quotes in verbatim environments
\IfFileExists{upquote.sty}{\usepackage{upquote}}{}
\IfFileExists{microtype.sty}{% use microtype if available
  \usepackage[]{microtype}
  \UseMicrotypeSet[protrusion]{basicmath} % disable protrusion for tt fonts
}{}
\makeatletter
\@ifundefined{KOMAClassName}{% if non-KOMA class
  \IfFileExists{parskip.sty}{%
    \usepackage{parskip}
  }{% else
    \setlength{\parindent}{0pt}
    \setlength{\parskip}{6pt plus 2pt minus 1pt}}
}{% if KOMA class
  \KOMAoptions{parskip=half}}
\makeatother
\usepackage{xcolor}
\usepackage{color}
\usepackage{fancyvrb}
\newcommand{\VerbBar}{|}
\newcommand{\VERB}{\Verb[commandchars=\\\{\}]}
\DefineVerbatimEnvironment{Highlighting}{Verbatim}{commandchars=\\\{\}}
% Add ',fontsize=\small' for more characters per line
\newenvironment{Shaded}{}{}
\newcommand{\AlertTok}[1]{\textcolor[rgb]{1.00,0.00,0.00}{\textbf{#1}}}
\newcommand{\AnnotationTok}[1]{\textcolor[rgb]{0.38,0.63,0.69}{\textbf{\textit{#1}}}}
\newcommand{\AttributeTok}[1]{\textcolor[rgb]{0.49,0.56,0.16}{#1}}
\newcommand{\BaseNTok}[1]{\textcolor[rgb]{0.25,0.63,0.44}{#1}}
\newcommand{\BuiltInTok}[1]{\textcolor[rgb]{0.00,0.50,0.00}{#1}}
\newcommand{\CharTok}[1]{\textcolor[rgb]{0.25,0.44,0.63}{#1}}
\newcommand{\CommentTok}[1]{\textcolor[rgb]{0.38,0.63,0.69}{\textit{#1}}}
\newcommand{\CommentVarTok}[1]{\textcolor[rgb]{0.38,0.63,0.69}{\textbf{\textit{#1}}}}
\newcommand{\ConstantTok}[1]{\textcolor[rgb]{0.53,0.00,0.00}{#1}}
\newcommand{\ControlFlowTok}[1]{\textcolor[rgb]{0.00,0.44,0.13}{\textbf{#1}}}
\newcommand{\DataTypeTok}[1]{\textcolor[rgb]{0.56,0.13,0.00}{#1}}
\newcommand{\DecValTok}[1]{\textcolor[rgb]{0.25,0.63,0.44}{#1}}
\newcommand{\DocumentationTok}[1]{\textcolor[rgb]{0.73,0.13,0.13}{\textit{#1}}}
\newcommand{\ErrorTok}[1]{\textcolor[rgb]{1.00,0.00,0.00}{\textbf{#1}}}
\newcommand{\ExtensionTok}[1]{#1}
\newcommand{\FloatTok}[1]{\textcolor[rgb]{0.25,0.63,0.44}{#1}}
\newcommand{\FunctionTok}[1]{\textcolor[rgb]{0.02,0.16,0.49}{#1}}
\newcommand{\ImportTok}[1]{\textcolor[rgb]{0.00,0.50,0.00}{\textbf{#1}}}
\newcommand{\InformationTok}[1]{\textcolor[rgb]{0.38,0.63,0.69}{\textbf{\textit{#1}}}}
\newcommand{\KeywordTok}[1]{\textcolor[rgb]{0.00,0.44,0.13}{\textbf{#1}}}
\newcommand{\NormalTok}[1]{#1}
\newcommand{\OperatorTok}[1]{\textcolor[rgb]{0.40,0.40,0.40}{#1}}
\newcommand{\OtherTok}[1]{\textcolor[rgb]{0.00,0.44,0.13}{#1}}
\newcommand{\PreprocessorTok}[1]{\textcolor[rgb]{0.74,0.48,0.00}{#1}}
\newcommand{\RegionMarkerTok}[1]{#1}
\newcommand{\SpecialCharTok}[1]{\textcolor[rgb]{0.25,0.44,0.63}{#1}}
\newcommand{\SpecialStringTok}[1]{\textcolor[rgb]{0.73,0.40,0.53}{#1}}
\newcommand{\StringTok}[1]{\textcolor[rgb]{0.25,0.44,0.63}{#1}}
\newcommand{\VariableTok}[1]{\textcolor[rgb]{0.10,0.09,0.49}{#1}}
\newcommand{\VerbatimStringTok}[1]{\textcolor[rgb]{0.25,0.44,0.63}{#1}}
\newcommand{\WarningTok}[1]{\textcolor[rgb]{0.38,0.63,0.69}{\textbf{\textit{#1}}}}
\usepackage{graphicx}
\makeatletter
\def\maxwidth{\ifdim\Gin@nat@width>\linewidth\linewidth\else\Gin@nat@width\fi}
\def\maxheight{\ifdim\Gin@nat@height>\textheight\textheight\else\Gin@nat@height\fi}
\makeatother
% Scale images if necessary, so that they will not overflow the page
% margins by default, and it is still possible to overwrite the defaults
% using explicit options in \includegraphics[width, height, ...]{}
\setkeys{Gin}{width=\maxwidth,height=\maxheight,keepaspectratio}
% Set default figure placement to htbp
\makeatletter
\def\fps@figure{htbp}
\makeatother
\setlength{\emergencystretch}{3em} % prevent overfull lines
\providecommand{\tightlist}{%
  \setlength{\itemsep}{0pt}\setlength{\parskip}{0pt}}
\setcounter{secnumdepth}{-\maxdimen} % remove section numbering
\ifLuaTeX
  \usepackage{selnolig}  % disable illegal ligatures
\fi
\IfFileExists{bookmark.sty}{\usepackage{bookmark}}{\usepackage{hyperref}}
\IfFileExists{xurl.sty}{\usepackage{xurl}}{} % add URL line breaks if available
\urlstyle{same} % disable monospaced font for URLs
\hypersetup{
  hidelinks,
  pdfcreator={LaTeX via pandoc}}

\author{}
\date{}

\begin{document}

\hypertarget{numerical-analysis-exercise-4}{%
\section{Numerical Analysis Exercise
4}\label{numerical-analysis-exercise-4}}

\hypertarget{4use-the-gaussian-elimination-with-backward-substitution-to-solve-the-following-liner-systems}{%
\subparagraph{4.Use the Gaussian elimination with backward substitution
to solve the following liner
systems.}\label{4use-the-gaussian-elimination-with-backward-substitution-to-solve-the-following-liner-systems}}

\[(1)\left\{\begin{matrix} 
  4x_{1}-x_{2}+3x_{3}=8 \\  
  2x_{1}+5x_{2}+2x_{3}=3.\\
  x_{1}+2x_{2}+4x_{3}=11
\end{matrix}\right.\]

\textbf{Solution}

\begin{Shaded}
\begin{Highlighting}[]
\VariableTok{clc}
\VariableTok{clear}
\CommentTok{\%Gaussian Elimination Method\%}
\VariableTok{A}\OperatorTok{=}\VariableTok{input}\NormalTok{(}\StringTok{"Please input the a:"}\NormalTok{)}
\VariableTok{b}\OperatorTok{=}\VariableTok{input}\NormalTok{(}\StringTok{"Please input the b:"}\NormalTok{)}
\VariableTok{b}\OperatorTok{=}\VariableTok{b}\OperatorTok{\textquotesingle{};}
\VariableTok{Gauss}\NormalTok{(}\VariableTok{A}\OperatorTok{,}\VariableTok{b}\NormalTok{)}
\end{Highlighting}
\end{Shaded}

\begin{Shaded}
\begin{Highlighting}[]
\KeywordTok{function} \VariableTok{X}\OperatorTok{=}\VariableTok{Gauss}\NormalTok{(}\VariableTok{A}\OperatorTok{,}\VariableTok{b}\NormalTok{)}
\VariableTok{ab}\OperatorTok{=}\NormalTok{[}\VariableTok{A}\OperatorTok{,}\VariableTok{b}\NormalTok{]}\OperatorTok{;}\CommentTok{\%a为系数矩阵，b为右端项}
\NormalTok{[}\VariableTok{R}\OperatorTok{,}\VariableTok{C}\NormalTok{]}\OperatorTok{=}\VariableTok{size}\NormalTok{(}\VariableTok{ab}\NormalTok{)}\OperatorTok{;}
\KeywordTok{for} \VariableTok{i}\OperatorTok{=}\FloatTok{1}\OperatorTok{:}\VariableTok{R}\OperatorTok{{-}}\FloatTok{1} \CommentTok{\%控制消去首行元素的次数 二阶方程组只需要消去一次,三阶消去两次                   i=2时，系数矩阵第三行消去一个首行元素,此时系数矩阵第三行只有一个元素}
    \KeywordTok{for} \VariableTok{j}\OperatorTok{=}\VariableTok{i}\OperatorTok{+}\FloatTok{1}\OperatorTok{:}\VariableTok{R}                                                                          \CommentTok{\%假设R=3，i=1时，j=2.此时内循环两次，系数矩阵的第二三行分别消去一个首行非零元素}
        \VariableTok{ab}\NormalTok{(}\VariableTok{j}\OperatorTok{,}\VariableTok{i}\OperatorTok{:}\VariableTok{C}\NormalTok{)}\OperatorTok{=}\VariableTok{ab}\NormalTok{(}\VariableTok{j}\OperatorTok{,}\VariableTok{i}\OperatorTok{:}\VariableTok{C}\NormalTok{)}\OperatorTok{{-}}\VariableTok{ab}\NormalTok{(}\VariableTok{i}\OperatorTok{,}\VariableTok{i}\OperatorTok{:}\VariableTok{C}\NormalTok{)}\OperatorTok{*}\VariableTok{ab}\NormalTok{(}\VariableTok{j}\OperatorTok{,}\VariableTok{i}\NormalTok{)}\OperatorTok{/}\VariableTok{ab}\NormalTok{(}\VariableTok{i}\OperatorTok{,}\VariableTok{i}\NormalTok{)}\OperatorTok{;}\CommentTok{\%每一行消去首行元素}
    \KeywordTok{end}
\KeywordTok{end}
\VariableTok{x}\OperatorTok{=}\VariableTok{zeros}\NormalTok{(}\VariableTok{R}\OperatorTok{,}\FloatTok{1}\NormalTok{)}\OperatorTok{;}\CommentTok{\%解空间}
\VariableTok{x}\NormalTok{(}\VariableTok{R}\OperatorTok{,}\FloatTok{1}\NormalTok{)}\OperatorTok{=}\VariableTok{ab}\NormalTok{(}\VariableTok{R}\OperatorTok{,}\VariableTok{C}\NormalTok{)}\OperatorTok{/}\VariableTok{ab}\NormalTok{(}\VariableTok{R}\OperatorTok{,}\VariableTok{R}\NormalTok{)}\OperatorTok{;}

\VariableTok{fprintf}\NormalTok{(}\StringTok{"解为："}\NormalTok{)}
\KeywordTok{for} \VariableTok{k}\OperatorTok{=}\VariableTok{R}\OperatorTok{{-}}\FloatTok{1}\OperatorTok{:{-}}\FloatTok{1}\OperatorTok{:}\FloatTok{1}\CommentTok{\%假设R=3，C=4；k从2开始循环}
    \VariableTok{x}\NormalTok{(}\VariableTok{k}\NormalTok{)}\OperatorTok{=}\NormalTok{(}\VariableTok{ab}\NormalTok{(}\VariableTok{k}\OperatorTok{,}\VariableTok{C}\NormalTok{)}\OperatorTok{{-}}\VariableTok{ab}\NormalTok{(}\VariableTok{k}\OperatorTok{,}\VariableTok{k}\OperatorTok{+}\FloatTok{1}\OperatorTok{:}\VariableTok{R}\NormalTok{)}\OperatorTok{*}\VariableTok{x}\NormalTok{(}\VariableTok{k}\OperatorTok{+}\FloatTok{1}\OperatorTok{:}\VariableTok{R}\NormalTok{))}\OperatorTok{/}\VariableTok{ab}\NormalTok{(}\VariableTok{k}\OperatorTok{,}\VariableTok{k}\NormalTok{)}\OperatorTok{;}\CommentTok{\%反向替换}
\KeywordTok{end}
\VariableTok{x}
\KeywordTok{end}
\end{Highlighting}
\end{Shaded}

\begin{Shaded}
\begin{Highlighting}[]
\VariableTok{A} \OperatorTok{=} \FloatTok{3}\NormalTok{×}\FloatTok{3}    
     \FloatTok{4}    \OperatorTok{{-}}\FloatTok{1}     \FloatTok{3}
     \FloatTok{2}     \FloatTok{5}     \FloatTok{2}
     \FloatTok{1}     \FloatTok{2}     \FloatTok{4}

\VariableTok{b} \OperatorTok{=} \FloatTok{1}\NormalTok{×}\FloatTok{3}    
     \FloatTok{8}     \FloatTok{3}    \FloatTok{11}

\VariableTok{R} \OperatorTok{=} \FloatTok{3}
\VariableTok{C} \OperatorTok{=} \FloatTok{4}
\NormalTok{解为：}
\VariableTok{x} \OperatorTok{=} \FloatTok{3}\NormalTok{×}\FloatTok{1}    
   \OperatorTok{{-}}\FloatTok{0.4328}
   \OperatorTok{{-}}\FloatTok{0.4627}
    \FloatTok{3.0896}
\end{Highlighting}
\end{Shaded}

\[(2)\left\{\begin{matrix} 
  4x_{1}+x_{2}+3x_{3}=9 \\  
  2x_{1}+4x_{2}-x_{3}=-5.\\
  x_{1}+2x_{2}-3x_{3}=-9
\end{matrix}\right.\]

\begin{Shaded}
\begin{Highlighting}[]
\VariableTok{A} \OperatorTok{=} \FloatTok{3}\NormalTok{×}\FloatTok{3}    
     \FloatTok{4}     \FloatTok{1}     \FloatTok{2}
     \FloatTok{2}     \FloatTok{4}    \OperatorTok{{-}}\FloatTok{1}
     \FloatTok{1}     \FloatTok{1}    \OperatorTok{{-}}\FloatTok{3}

\VariableTok{b} \OperatorTok{=} \FloatTok{1}\NormalTok{×}\FloatTok{3}    
     \FloatTok{9}    \OperatorTok{{-}}\FloatTok{5}    \OperatorTok{{-}}\FloatTok{9}

\NormalTok{解为：}
\VariableTok{x} \OperatorTok{=} \FloatTok{3}\NormalTok{×}\FloatTok{1}    
     \FloatTok{1}
    \OperatorTok{{-}}\FloatTok{1}
     \FloatTok{3}
\end{Highlighting}
\end{Shaded}

\hypertarget{5find-the-parabola-uxayuxauxaauxaxuxa2uxauxabuxaxuxauxacuxa-that-passes-throgh-the-points-1427-and-314}{%
\subparagraph{\texorpdfstring{5.Find the parabola \(y=Ax^{2}+Bx+C\) that
passes throgh the points (1,4),(2,7) and
(3,14).}{5.Find the parabola y=Ax\^{}\{2\}+Bx+C that passes throgh the points (1,4),(2,7) and (3,14).}}\label{5find-the-parabola-uxayuxauxaauxaxuxa2uxauxabuxaxuxauxacuxa-that-passes-throgh-the-points-1427-and-314}}

\textbf{Solution} Since the curve passer throgh these points,so we plug
the points into the equation,

\[\left\{\begin{matrix} 
  A+B+C=4 \\  
  4A+2B+C=7.\\
  9A+3B+C=14
\end{matrix}\right.\]

By using the model in the fourth question,we have,

\begin{Shaded}
\begin{Highlighting}[]
\VariableTok{A} \OperatorTok{=} \FloatTok{3}\NormalTok{×}\FloatTok{3}    
     \FloatTok{1}     \FloatTok{1}     \FloatTok{1}
     \FloatTok{4}     \FloatTok{2}     \FloatTok{1}
     \FloatTok{9}     \FloatTok{3}     \FloatTok{1}

\VariableTok{b} \OperatorTok{=} \FloatTok{1}\NormalTok{×}\FloatTok{3}    
     \FloatTok{4}     \FloatTok{7}    \FloatTok{14}

\NormalTok{解为：}
\VariableTok{x} \OperatorTok{=} \FloatTok{3}\NormalTok{×}\FloatTok{1}    
     \FloatTok{2}
    \OperatorTok{{-}}\FloatTok{3}
     \FloatTok{5}
\end{Highlighting}
\end{Shaded}

So ,we have the A = 2, B = -3 , C =5 . thus, we use the matlab figuring

\begin{figure}
\centering
\includegraphics{G:/Numercial Analysis Exercise/作业/Numerical Analysis Exercise 4/fig.png}
\caption{}
\end{figure}

\end{document}
